\section{Bài tập chuyên đề}

Hệ thống bài tập được sắp xếp theo mức độ tăng dần. Phần năng lượng chiếm trọng tâm của chuyên đề; phần tốc độ phản ứng chỉ giữ các dạng phù hợp với yêu cầu lớp 8 và một số bài đọc dữ liệu cơ bản.

\subsection{Củng cố}

\begin{tuluyen}
Phân loại các quá trình sau thành tỏa nhiệt hoặc thu nhiệt: (a) đốt than; (b) nung \ce{CaCO3}; (c) đốt ethanol; (d) cho \ce{CaO} vào nước.
\dapso{Tỏa nhiệt: (a), (c), (d); thu nhiệt: (b).}
\begin{loigiaiBT}
Các phản ứng cháy và phản ứng \ce{CaO + H2O -> Ca(OH)2} giải phóng nhiệt; nhiệt phân \ce{CaCO3} cần cung cấp nhiệt.
\textbf{Vậy,} (a), (c), (d) là tỏa nhiệt; (b) là thu nhiệt.
\end{loigiaiBT}
\end{tuluyen}

\begin{tuluyen}
Một loại than có nhiệt trị \SI{30}{\mega\joule\per\kilo\gram}. Tính nhiệt lượng tỏa ra khi đốt hoàn toàn \SI{12}{\kilo\gram} than.
\dapso{\SI{360}{\mega\joule}.}
\begin{loigiaiBT}
\[
Q=12\times30=\SI{360}{\mega\joule}.
\]
\textbf{Vậy,} nhiệt lượng tỏa ra là \SI{360}{\mega\joule}.
\end{loigiaiBT}
\end{tuluyen}

\begin{tuluyen}
Khi đốt hoàn toàn \(1\) mol \ce{H2}, phản ứng tỏa \SI{286}{\kilo\joule}. Tính nhiệt lượng tỏa ra khi đốt \SI{2.0}{\kilo\gram} \ce{H2}.
\dapso{\SI{286}{\mega\joule}.}
\begin{loigiaiBT}
\[
n_{\ce{H2}}=\frac{2000}{2}=\SI{1000}{\mole},
\]
\[
Q=1000\times286=\SI{286000}{\kilo\joule}=\SI{286}{\mega\joule}.
\]
\textbf{Vậy,} nhiệt lượng tỏa ra là \SI{286}{\mega\joule}.
\end{loigiaiBT}
\end{tuluyen}

\begin{tuluyen}
Một thiết bị cần \SI{72}{\mega\joule} năng lượng có ích. Hiệu suất sử dụng nhiệt là \(60\%\). Tính nhiệt lượng nhiên liệu phải tỏa ra.
\dapso{\SI{120}{\mega\joule}.}
\begin{loigiaiBT}
\[
Q_{\text{tỏa}}=\frac{72}{0.60}=\SI{120}{\mega\joule}.
\]
\textbf{Vậy,} nhiên liệu phải tỏa ra \SI{120}{\mega\joule}.
\end{loigiaiBT}
\end{tuluyen}

\begin{tuluyen}
Một nhiên liệu có nhiệt trị \SI{40}{\mega\joule\per\kilo\gram}. Bếp có hiệu suất \(50\%\). Tính khối lượng nhiên liệu cần dùng để cung cấp \SI{100}{\mega\joule} nhiệt có ích.
\dapso{\SI{5.0}{\kilo\gram}.}
\begin{loigiaiBT}
Mỗi kilogram nhiên liệu cho năng lượng có ích:
\[
40\times0.50=\SI{20}{\mega\joule}.
\]
Do đó:
\[
m=\frac{100}{20}=\SI{5.0}{\kilo\gram}.
\]
\textbf{Vậy,} cần \SI{5.0}{\kilo\gram} nhiên liệu.
\end{loigiaiBT}
\end{tuluyen}

\subsection{Vận dụng}

\begin{tuluyen}
LPG gồm \(50\%\) propane và \(50\%\) butane theo số mol. Khi đốt hoàn toàn, \(1\) mol propane và \(1\) mol butane lần lượt tỏa \SI{2220}{\kilo\joule} và \SI{2850}{\kilo\joule}. Tính nhiệt lượng tỏa ra khi đốt \(80\) mol hỗn hợp.
\source{Phỏng theo dạng bài LPG trong đề tuyển sinh chuyên Hóa Điện Biên 2026.}
\dapso{\SI{202.8}{\mega\joule}.}
\begin{loigiaiBT}
Có \(40\) mol mỗi khí:
\[
Q=40\times2220+40\times2850=\SI{202800}{\kilo\joule}
=\SI{202.8}{\mega\joule}.
\]
\textbf{Vậy,} nhiệt lượng tỏa ra là \SI{202.8}{\mega\joule}.
\end{loigiaiBT}
\end{tuluyen}

\begin{tuluyen}
Một bình LPG chứa \SI{12}{\kilo\gram} nhiên liệu có nhiệt trị trung bình \SI{46}{\mega\joule\per\kilo\gram}. Bếp sử dụng có hiệu suất \(65\%\). Gia đình cần trung bình \SI{15}{\mega\joule} nhiệt có ích mỗi ngày. Ước tính số ngày dùng hết bình gas.
\source{Phỏng theo dạng bài LPG và chi phí sinh hoạt trong đề tuyển sinh chuyên Hóa Lâm Đồng 2026.}
\dapso{\(\approx23.9\) ngày.}
\begin{loigiaiBT}
Tổng nhiệt lượng tỏa ra:
\[
Q=12\times46=\SI{552}{\mega\joule}.
\]
Nhiệt có ích:
\[
Q_{\text{ích}}=552\times0.65=\SI{358.8}{\mega\joule}.
\]
Số ngày:
\[
t=\frac{358.8}{15}\approx23.9.
\]
\textbf{Vậy,} một bình gas dùng được khoảng \(23.9\) ngày.
\end{loigiaiBT}
\end{tuluyen}

\begin{tuluyen}
Một xe dùng \SI{5.0}{\liter} nhiên liệu có khối lượng riêng \SI{0.72}{\kilo\gram\per\liter}, nhiệt trị \SI{44}{\mega\joule\per\kilo\gram}. Hiệu suất động cơ là \(28\%\), xe cần \SI{550}{\kilo\joule} công cơ học cho mỗi kilomet. Tính quãng đường tối đa xe đi được.
\source{Phỏng theo dạng bài E5/E10 trong đề tuyển sinh chuyên Hóa Hà Nội 2025 và Lâm Đồng 2026.}
\dapso{\(\approx80.6\) km.}
\begin{loigiaiBT}
Khối lượng nhiên liệu:
\[
m=5.0\times0.72=\SI{3.6}{\kilo\gram}.
\]
Nhiệt lượng tỏa ra:
\[
Q=3.6\times44=\SI{158.4}{\mega\joule}=\SI{158400}{\kilo\joule}.
\]
Công có ích:
\[
Q_{\text{ích}}=158400\times0.28=\SI{44352}{\kilo\joule}.
\]
Quãng đường:
\[
s=\frac{44352}{550}\approx80.6\text{ km}.
\]
\textbf{Vậy,} xe đi được tối đa khoảng \(80.6\) km.
\end{loigiaiBT}
\end{tuluyen}

\begin{tuluyen}
Một loại gas giá \(32\,000\) đồng/kg, nhiệt trị \SI{46}{\mega\joule\per\kilo\gram}. Bếp có hiệu suất \(60\%\). Tính chi phí để cung cấp \SI{414}{\mega\joule} nhiệt có ích.
\dapso{\(480\,000\) đồng.}
\begin{loigiaiBT}
Mỗi kilogram gas cung cấp:
\[
46\times0.60=\SI{27.6}{\mega\joule}
\]
nhiệt có ích. Khối lượng gas:
\[
m=\frac{414}{27.6}=\SI{15}{\kilo\gram}.
\]
Chi phí:
\[
C=15\times32000=480000\text{ đồng}.
\]
\textbf{Vậy,} chi phí là \(480\,000\) đồng.
\end{loigiaiBT}
\end{tuluyen}

\begin{tuluyen}
Để nung \SI{1.00}{\ton} \ce{CaCO3} cần cung cấp \SI{178}{\kilo\joule} cho mỗi mol \ce{CaCO3}. Lò dùng than có nhiệt trị \SI{25}{\mega\joule\per\kilo\gram}, hiệu suất truyền nhiệt \(40\%\). Tính khối lượng than tối thiểu. Cho \(M_{\ce{CaCO3}}=100\) g/mol.
\dapso{\SI{178}{\kilo\gram}.}
\begin{loigiaiBT}
\[
n_{\ce{CaCO3}}=\frac{1000000}{100}=\SI{10000}{\mole}.
\]
Năng lượng cần cho phản ứng:
\[
Q_{\text{ích}}=10000\times178=\SI{1780}{\mega\joule}.
\]
Than phải tỏa:
\[
Q=\frac{1780}{0.40}=\SI{4450}{\mega\joule}.
\]
Khối lượng than:
\[
m=\frac{4450}{25}=\SI{178}{\kilo\gram}.
\]
\textbf{Vậy,} cần tối thiểu \SI{178}{\kilo\gram} than.
\end{loigiaiBT}
\end{tuluyen}

\begin{tuluyen}
Methane cháy theo phương trình:
\[
\ce{CH4 + 2O2 -> CO2 + 2H2O}.
\]
Biết \(1\) mol methane tỏa \SI{890}{\kilo\joule}. Tính khối lượng \ce{CO2} tạo thành khi methane cung cấp \SI{8900}{\kilo\joule}, giả sử cháy hoàn toàn.
\dapso{\SI{440}{\gram}.}
\begin{loigiaiBT}
Số mol methane:
\[
n_{\ce{CH4}}=\frac{8900}{890}=\SI{10}{\mole}.
\]
Theo phương trình:
\[
n_{\ce{CO2}}=\SI{10}{\mole},
\]
nên:
\[
m_{\ce{CO2}}=10\times44=\SI{440}{\gram}.
\]
\textbf{Vậy,} khối lượng \ce{CO2} tạo thành là \SI{440}{\gram}.
\end{loigiaiBT}
\end{tuluyen}

\subsection{Tổng hợp}

\begin{tuluyen}
Một nhiên liệu lỏng gồm \(10\%\) ethanol theo thể tích, phần còn lại được xem như hydrocarbon X. Cho \SI{4.0}{\liter} nhiên liệu. Ethanol có \(D=\SI{0.789}{\kilo\gram\per\liter}\), X có \(D=\SI{0.72}{\kilo\gram\per\liter}\). Nhiệt trị theo khối lượng của ethanol và X lần lượt là \SI{27}{\mega\joule\per\kilo\gram} và \SI{44}{\mega\joule\per\kilo\gram}. Hiệu suất động cơ là \(30\%\); xe cần \SI{600}{\kilo\joule} công cơ học để đi \SI{1}{\kilo\metre}. Tính quãng đường tối đa.
\source{Phỏng theo cấu trúc bài nhiên liệu E10 trong đề tuyển sinh chuyên Hóa Lâm Đồng 2026.}
\dapso{\(\approx\SI{61.3}{\kilo\metre}\).}
\begin{loigiaiBT}
Thể tích ethanol là \SI{0.4}{\liter}, nên:
\[
m_{\text{ethanol}}=0.4\times0.789=\SI{0.3156}{\kilo\gram}.
\]
Thể tích X là \SI{3.6}{\liter}, nên:
\[
m_X=3.6\times0.72=\SI{2.592}{\kilo\gram}.
\]
Tổng nhiệt lượng:
\[
Q=0.3156\times27+2.592\times44
\approx\SI{122.57}{\mega\joule}.
\]
Công có ích:
\[
Q_{\text{ích}}=122.57\times0.30\approx\SI{36.77}{\mega\joule}
=\SI{36770}{\kilo\joule}.
\]
Quãng đường:
\[
s=\frac{36770}{600}\approx\SI{61.3}{\kilo\metre}.
\]
\textbf{Vậy,} xe đi được khoảng \SI{61.3}{\kilo\metre}.
\end{loigiaiBT}
\end{tuluyen}

\begin{tuluyen}
Hai nhiên liệu A và B có dữ kiện:
\begin{center}
\begin{tabular}{c|ccc}
\toprule
Nhiên liệu & Nhiệt trị (MJ/kg) & Hiệu suất & Giá (đồng/kg) \\
\midrule
A & 45 & 60\% & 30000 \\
B & 32 & 75\% & 21000 \\
\bottomrule
\end{tabular}
\end{center}
Tính chi phí để mỗi nhiên liệu cung cấp \SI{270}{\mega\joule} năng lượng có ích và cho biết phương án rẻ hơn.
\source{Phỏng theo các bài so sánh nhiên liệu và kinh tế trong đề chuyên Hóa Đắk Lắk 2026 và Lâm Đồng 2026.}
\dapso{A: \(300\,000\) đồng; B: \(236\,250\) đồng; B rẻ hơn.}
\begin{loigiaiBT}
A cho năng lượng có ích trên mỗi kilogram:
\[
45\times0.60=\SI{27}{\mega\joule\per\kilo\gram}.
\]
Cần:
\[
m_A=\frac{270}{27}=\SI{10}{\kilo\gram},
\]
chi phí:
\[
C_A=10\times30000=300000\text{ đồng}.
\]
B cho:
\[
32\times0.75=\SI{24}{\mega\joule\per\kilo\gram}.
\]
Cần:
\[
m_B=\frac{270}{24}=\SI{11.25}{\kilo\gram},
\]
chi phí:
\[
C_B=11.25\times21000=236250\text{ đồng}.
\]
\textbf{Vậy,} nhiên liệu B có chi phí thấp hơn.
\end{loigiaiBT}
\end{tuluyen}

\begin{tuluyen}
Một hỗn hợp khí gồm \(30\%\) methane và \(70\%\) propane theo số mol. \(1\) mol methane và propane khi cháy lần lượt tỏa \SI{890}{\kilo\joule} và \SI{2220}{\kilo\joule}. Một hệ thống cần \SI{500}{\mega\joule} năng lượng có ích, hiệu suất \(50\%\). Tính số mol hỗn hợp khí tối thiểu cần đốt.
\source{Phỏng theo kiểu bài nhiên liệu hỗn hợp trong đề chuyên Hóa Điện Biên 2026.}
\dapso{\(\approx\SI{549.1}{\mole}\).}
\begin{loigiaiBT}
Nhiệt lượng hỗn hợp phải tỏa:
\[
Q=\frac{500}{0.50}=\SI{1000}{\mega\joule}=\SI{1000000}{\kilo\joule}.
\]
Một mol hỗn hợp tỏa:
\[
q=0.30\times890+0.70\times2220=\SI{1821}{\kilo\joule\per\mole}.
\]
Số mol hỗn hợp:
\[
n=\frac{1000000}{1821}\approx\SI{549.1}{\mole}.
\]
\textbf{Vậy,} cần khoảng \SI{549.1}{\mole} hỗn hợp khí.
\end{loigiaiBT}
\end{tuluyen}

\begin{tuluyen}
Hydrogen và methane đều được dùng làm nhiên liệu. Biết:
\[
\ce{2H2 + O2 -> 2H2O},\qquad
\ce{CH4 + 2O2 -> CO2 + 2H2O}.
\]
\(1\) mol \ce{H2} tỏa \SI{286}{\kilo\joule}, \(1\) mol \ce{CH4} tỏa \SI{890}{\kilo\joule}. Một xe cần \SI{89}{\mega\joule} năng lượng hóa học. Tính khối lượng mỗi nhiên liệu cần dùng và lượng \ce{CO2} tạo ra.
\source{Phỏng theo dạng bài so sánh gasoline/LPG/\ce{H2} trong đề chuyên Hóa Đắk Lắk 2026.}
\dapso{\(\approx\SI{0.622}{\kilo\gram}\) \ce{H2}, không tạo \ce{CO2}; \(\SI{1.60}{\kilo\gram}\) \ce{CH4}, tạo \(\SI{4.40}{\kilo\gram}\) \ce{CO2}.}
\begin{loigiaiBT}
Với hydrogen:
\[
n_{\ce{H2}}=\frac{89000}{286}\approx\SI{311.2}{\mole},
\]
\[
m_{\ce{H2}}\approx311.2\times2=\SI{622.4}{\gram}.
\]
Phản ứng không tạo \ce{CO2}.

Với methane:
\[
n_{\ce{CH4}}=\frac{89000}{890}=\SI{100}{\mole},
\]
\[
m_{\ce{CH4}}=100\times16=\SI{1.60}{\kilo\gram}.
\]
Theo phương trình, tạo \(100\) mol \ce{CO2}:
\[
m_{\ce{CO2}}=100\times44=\SI{4.40}{\kilo\gram}.
\]
\textbf{Vậy,} để cung cấp cùng \SI{89}{\mega\joule}, cần khoảng \SI{0.622}{\kilo\gram} \ce{H2} hoặc \SI{1.60}{\kilo\gram} \ce{CH4}; methane tạo \SI{4.40}{\kilo\gram} \ce{CO2}.
\end{loigiaiBT}
\end{tuluyen}

\begin{tuluyen}
Trong một quy trình, để sản xuất một mẻ sản phẩm cần \SI{300}{\mega\joule} nhiệt có ích. Lò có hiệu suất \(60\%\). Có hai lựa chọn:
\begin{itemize}
\item nhiên liệu A: nhiệt trị \SI{50}{\mega\joule\per\kilo\gram}, giá \(35\,000\) đồng/kg;
\item nhiên liệu B: nhiệt trị \SI{36}{\mega\joule\per\kilo\gram}, giá \(22\,000\) đồng/kg.
\end{itemize}
Tính khối lượng và chi phí cho mỗi phương án, từ đó lựa chọn theo tiêu chí chi phí.
\dapso{A: \SI{10}{\kilo\gram}, \(350\,000\) đồng; B: \(\approx\SI{13.89}{\kilo\gram}\), \(\approx305\,600\) đồng; chọn B nếu chỉ xét chi phí.}
\begin{loigiaiBT}
Nhiệt lượng nhiên liệu phải tỏa:
\[
Q=\frac{300}{0.60}=\SI{500}{\mega\joule}.
\]
Với A:
\[
m_A=\frac{500}{50}=\SI{10}{\kilo\gram},
\qquad C_A=10\times35000=350000\text{ đồng}.
\]
Với B:
\[
m_B=\frac{500}{36}\approx\SI{13.89}{\kilo\gram},
\]
\[
C_B\approx13.89\times22000\approx305600\text{ đồng}.
\]
\textbf{Vậy,} nếu chỉ xét chi phí nhiên liệu, phương án B kinh tế hơn.
\end{loigiaiBT}
\end{tuluyen}

\subsection{Tốc độ phản ứng}

\begin{tuluyen}
Hai cốc chứa cùng khối lượng \ce{CaCO3} và cùng thể tích \ce{HCl} dư. Cốc A dùng đá vôi dạng cục, cốc B dùng đá vôi nghiền nhỏ. Các điều kiện khác như nhau. Cốc nào tạo \ce{CO2} nhanh hơn?
\dapso{Cốc B.}
\begin{loigiaiBT}
Đá vôi nghiền nhỏ có tổng diện tích bề mặt tiếp xúc với acid lớn hơn nên phản ứng nhanh hơn.
\textbf{Vậy,} cốc B tạo \ce{CO2} nhanh hơn.
\end{loigiaiBT}
\end{tuluyen}

\begin{tuluyen}
Cùng một mẩu Zn được cho vào hai dung dịch \ce{HCl} có cùng thể tích, một dung dịch \SI{0.5}{\mole\per\liter}, một dung dịch \SI{1.5}{\mole\per\liter}; acid đều dư. Dung dịch nào làm Zn phản ứng nhanh hơn? Lượng \ce{H2} cuối cùng có khác nhau không?
\dapso{Dung dịch \SI{1.5}{\mole\per\liter} làm phản ứng nhanh hơn; lượng \ce{H2} cuối cùng bằng nhau.}
\begin{loigiaiBT}
Nồng độ acid lớn hơn làm phản ứng nhanh hơn. Do lượng Zn ban đầu như nhau và acid đều dư, lượng \ce{H2} cuối cùng do Zn quyết định.
\textbf{Vậy,} dung dịch \SI{1.5}{\mole\per\liter} cho tốc độ lớn hơn nhưng lượng \ce{H2} cuối cùng không đổi.
\end{loigiaiBT}
\end{tuluyen}

\begin{tuluyen}
Một phản ứng tạo khí cho bảng:
\[
\begin{array}{c|cccc}
t\text{ (s)} & 0 & 20 & 40 & 60\\ \hline
V\text{ (mL)} & 0 & 38 & 64 & 76
\end{array}
\]
Trong hai khoảng 0--20 s và 40--60 s, khoảng nào phản ứng nhanh hơn? Giải thích bằng mức tăng thể tích khí trung bình.
\source{Phỏng theo dạng đọc dữ liệu tốc độ trong đề chuyên Hóa Điện Biên và Đà Nẵng.}
\dapso{0--20 s nhanh hơn; lần lượt \SI{1.9}{\milli\liter\per\second} và \SI{0.6}{\milli\liter\per\second}.}
\begin{loigiaiBT}
\[
\overline v_{0-20}=\frac{38}{20}=\SI{1.9}{\milli\liter\per\second},
\]
\[
\overline v_{40-60}=\frac{76-64}{20}=\SI{0.6}{\milli\liter\per\second}.
\]
\textbf{Vậy,} phản ứng nhanh hơn trong khoảng 0--20 s.
\end{loigiaiBT}
\end{tuluyen}

\begin{tuluyen}
Trong phản ứng phân hủy \ce{H2O2}, thêm \ce{MnO2} làm oxygen thoát ra nhanh hơn. Nêu vai trò của \ce{MnO2} và cho biết chất này có làm tăng lượng oxygen lí thuyết thu được từ cùng lượng \ce{H2O2} hay không.
\dapso{\ce{MnO2} là chất xúc tác; không làm tăng lượng \ce{O2} lí thuyết từ cùng lượng \ce{H2O2}.}
\begin{loigiaiBT}
\ce{MnO2} làm phản ứng diễn ra nhanh hơn nhưng không bị tiêu hao theo phương trình tổng quát. Lượng \ce{O2} lí thuyết vẫn do lượng \ce{H2O2} quyết định.
\textbf{Vậy,} \ce{MnO2} là chất xúc tác và không làm tăng lượng \ce{O2} lí thuyết.
\end{loigiaiBT}
\end{tuluyen}
